\begin{center}
  \large\textbf{ABSTRACT}
\end{center}

\addcontentsline{toc}{chapter}{ABSTRACT}

\vspace{2ex}

\begingroup
% Menghilangkan padding
\setlength{\tabcolsep}{0pt}

\noindent
\begin{tabularx}{\textwidth}{l >{\centering}m{3em} X}
  \emph{Name}     & : & Yoel Yhokhanan Sianipar         \\

  \emph{Title}    & : & \emph{INVERTER-SUPPLIED WATER PUMP PRESSURE CONTROL SYSTEM USING NEURAL NETWORK}   \\

  \emph{Advisors} & : & 1. Prof. Dr. Muhammad Rivai, S.T., M.T.   \\
                  &   & 2. Dr. Suwito, S.T., M.T. \\
\vspace{1ex}
\end{tabularx}
\endgroup

% Ubah paragraf berikut dengan abstrak dari tugas akhir dalam Bahasa Inggris
\emph{Unstable water pressure in household water pumps often leads to inefficiencies in water distribution, 
especially when multiple faucets are opened simultaneously. The current solution involves using additional 
external devices such as a pressure switch. This study offers an alternative solution by controlling the pump’s 
speed and pressure through input supply regulation using a PWM AC Chopper. The system is developed by incorporating 
Neural Network (NN) technology as the main controller, which adaptively adjusts the 220 VAC inverter output voltage 
based on varying load conditions. 
The Neural Network is trained to recognize water pressure patterns and adjust the output in real time, thereby 
maintaining stable pressure even under changing system loads. With this approach, the system is expected to deliver more responsive 
and precise performance. This research aims to contribute to the development of intelligent control systems for water pumps 
using neural network methods.}

% Ubah kata-kata berikut dengan kata kunci dari tugas akhir dalam Bahasa Inggris
\emph{Keywords}: \emph{PWM AC Chopper}, \emph{Neural Network}, \emph{Water Pump Pressure}.
